\documentclass[11pt,a4paper]{article}
\usepackage[utf8]{inputenc}\usepackage[T1]{fontenc}
\usepackage[margin=1in]{geometry}
\usepackage{amsmath,amssymb,booktabs,graphicx,enumitem,xcolor,parskip,titlesec}
\usepackage[hidelinks]{hyperref}
\definecolor{qnavy}{RGB}{20,40,75}\definecolor{qgrey}{RGB}{90,90,90}\definecolor{qred}{RGB}{150,30,30}
\definecolor{qgreen}{RGB}{47,107,79}
\titleformat{\section}{\large\bfseries\color{qnavy}}{\thesection}{0.6em}{}
\begin{document}\thispagestyle{empty}
\noindent
{\large\bfseries\color{qnavy} Report R19 --- Evidence for Two Numbers Paper 1 Had Asserted}\\[2pt]
{\color{qgrey}Paper 1's external site-probe rows had no persisted artefact, and the input file
named ``he'' was not H\&E. Both rows now have evidence and the H\&E restriction changes
nothing.}\\[6pt]
{\color{qgrey}\rule{\textwidth}{0.4pt}}

\medskip
\noindent 17 August 2026 \quad\textbullet\quad Quantara \quad\textbullet\quad
{\color{qgrey}Peer-review audit copy}

\section*{Summary}

Paper 1 reports three probes showing that site identity is recoverable from whole-slide features.
While preparing the second-encoder arm, two problems surfaced with the two most striking rows.

\textbf{First, they had no evidence file.} Only the within-TCGA row (0.7062) is backed by a
persisted artefact --- R12's \texttt{results.json}, field
\texttt{silo\_signature\_probe.balanced\_accuracy}. The external rows, 1.0000 and 0.9957, were
computed in a session and transcribed into the manuscript. Nothing was written. They were
unverifiable, and this programme's claim that Paper 1 was ``23/23 verified'' was asserted over
them.

\textbf{Second, the input was not what its name said.} The probe consumed
\texttt{external\_he\_mean.npz}. Despite \texttt{he} in the filename, 3 of its 29 HTMCP slides are
not H\&E --- \texttt{HTMCP-02-06-01032-CHR} and \texttt{HTMCP-02-15-01090-CHR} (chromogenic ISH)
and \texttt{HTMCP-02-07-01065-06-P16} (p16 immunohistochemistry). R12's audit had already flagged
that HTMCP-LC is majority IHC and that only its H\&E subset is a valid unit. The filter was then
not enforced in the file asserting it had been. With balanced accuracy at exactly 1.000, a stain
difference is the cheapest possible shortcut, so this was not a cosmetic problem.

\textbf{Both rows now reproduce, and the H\&E restriction costs nothing.}

\begin{center}
\begin{tabular}{@{}llrrrr@{}}
\toprule
& \textbf{Probe} & \textbf{Bal.\ acc.} & \textbf{Chance} & \textbf{Null mean} &
\textbf{Perm.\ $p$} \\
\midrule
ARM\_ALL & EAGLE vs HTMCP, published 78 slides & \textbf{1.0000} & 0.5000 & 0.5074 & 0.0050 \\
ARM\_HE & EAGLE vs HTMCP, \textbf{H\&E only}, 75 & \textbf{1.0000} & 0.5000 & 0.5058 & 0.0050 \\
\midrule
ARM3\_ALL & TCGA vs EAGLE vs HTMCP, balanced & 0.9816 & 0.3333 & 0.3345 & (see \S3) \\
ARM3\_HE & the same, \textbf{H\&E only} & 0.9814 & 0.3333 & 0.3660 & (see \S3) \\
\bottomrule
\end{tabular}
\end{center}

\noindent\textcolor{qgreen}{\textbf{The contribution of the three non-H\&E slides is exactly
0.0000.}} Verdict \textbf{EXTERNAL\_SITE\_SIGNATURE\_CONFIRMED\_HE\_ONLY}, from rules
pre-declared in the config and hashed at \texttt{5985b7c75d064a20\ldots} before any arm ran.

\section{What was and was not reproduced}

\begin{itemize}[leftmargin=1.3em,itemsep=3pt]
\item \textbf{Row 2 reproduces exactly.} 1.0000, and the null lands at 0.5074 against the
      manuscript's 0.5052 --- a difference consistent with a different permutation draw. The
      manuscript's number stands as written.
\item \textbf{Row 3 does not reproduce exactly:} 0.9816 here against 0.9957 published, a gap of
      0.014. \textbf{The original probe's hyperparameters were never recorded}, so there is no way
      to tell whether the difference is the model, the cross-validation, the number of subsampling
      draws, or the class balance. This run declares its own --- standardise, logistic regression
      $C = 1.0$, 5-fold stratified out-of-fold balanced accuracy, 20 size-balanced draws at 29 (or
      26) slides per cohort --- so that the number is at least reproducible from now on. Paper 1
      should quote 0.9816 with this method rather than 0.9957 with none.
\item \textbf{Row 1 was not recomputed.} It already has an artefact and was left alone.
\end{itemize}

The gate ordering matters here: G4 required the published slide set to recover 1.0000 \emph{before}
the H\&E arm was interpreted, because if the manuscript number were not reproducible from the
surviving inputs then that, and not the stain question, would be the result.

\section{Why the three-way $p$ is not quoted as 0.048}

The 3-way arms return $p = 0.0476$, which is the floor: with 20 subsampling draws the smallest
attainable value is $1/21$. That number is an artefact of the design, not a measurement, so it is
not reported as evidential. What is evidential is the separation --- 0.98 against a null centred at
0.3345 and 0.3660, with a standard deviation across draws of 0.013 and 0.019. Paper 1's
$p = 0.0050$ for this row came from a permutation scheme that is not recorded and cannot be
reconstructed; it should not be quoted either.

Both two-cohort arms do reach $p = 0.0050$ legitimately, from 200 label permutations under
$p = (1 + \#\{\text{null} \geq \text{obs}\})/(B+1)$.

\section{The honest reading of a perfect score}

Balanced accuracy of 1.000 on 75 slides is separation, not prediction. With 1024 features and 75
samples a linear probe can separate many arbitrary labellings, which is precisely why the
permutation null carries the weight: it sits at 0.5058 with a 95th percentile of 0.617, so the
observed 1.000 is far outside anything label noise produces here. That is the argument, and it
survives the stain restriction.

It remains unsurprising on the merits. An Italian population-based series and a sub-Saharan
HIV-positive series differ in fixation, stain lot, scanner and protocol, and a representation
trained without invariance objectives will encode that. The finding is that it is total, not that
it exists.

\section{Limitations}

\textbf{Stain was determined from slide-id suffixes.} \texttt{HE}, \texttt{CHR}, \texttt{P16} in
the filename are the only stain field available in the staged features. A slide mislabelled at
source would be misclassified here.

\textbf{75 slides, two institutions.} This bounds the site signature's existence across genuinely
external acquisition. It is not a powered external benchmark and R12's audit already said so.

\textbf{Mean-pooled slide features.} The probe is over slide-level means, not the attention-MIL
representation, so it measures what is linearly recoverable from average tile content.

\textbf{Row 3's discrepancy is unresolved and unresolvable.} Without the original
hyperparameters the 0.014 gap cannot be attributed. The remedy is forward-looking only.

\textbf{One encoder.} These are Phikon-v2 features. Whether the same holds for a second encoder is
R18, in flight at the time of writing.

\section*{Reproducing this report}

\texttt{evidence/m15\_external\_site\_probe.py} (the whole run),
\texttt{evidence/config.yaml} (pre-declared rules, SHA-256
\texttt{5985b7c75d064a20\ldots}), \texttt{evidence/gates.json} (6 gates),
\texttt{evidence/results.json}, \texttt{evidence/log.jsonl},
\texttt{evidence/probe\_set.csv} (all 78 slides with cohort and H\&E flag).
Input: \texttt{quantara-staging/ext\_feat/external\_he\_mean.npz} and
\texttt{armC/data/runs/mean\_embeddings.npz}, both hashed in
\texttt{evidence/inputs.json}.

\end{document}
